\section*{Scriptural Date and Location}

\begin{scripturestable}
\scripturerow{First Reading}{Exod.~16:2--4, 12--15}{The wilderness of Sin,
between Elim and Sinai}{Narrated in Israel's wilderness journey}
\scripturenote{Exodus presents Israel's murmuring and God's provision within
the journey from Egypt to Sinai. The book's final literary form addresses
Israel through remembered deliverance and covenant formation; the narrated
event and the composition of the received text are not the same date.}
\scripturerow{Entrance}{Ps.~69:2, 6 (modern 70)}{No place named in the appointed verses}{Before c.~165 BC; individual date uncertain}
\scripturenote{The Douay--Rheims title names David, but neither that attribution nor the appointed pleas for help establishes an occasion or composition year. The \sourceurl{https://bible.usccb.org/bible/psalms/0}{NABRE Psalms introduction} supplies only the collection's latest composition boundary and denies secure dating of an individual psalm.}
\scripturerow{Psalm}{Ps.~78:3--4, 23--25, 54}{Israel; the remembered route
runs from Egypt and wilderness to the sanctuary land}{A historical psalm
looking back over generations}
\scripturenote{The psalm speaks as instruction received from the ancestors and
retold to a later generation. Its rehearsal of manna interprets the
wilderness event inside a longer history of divine fidelity and human
forgetfulness.}
\scripturerow{Communion A}{Wis.~16:20}{Alexandria in Gigot's account}{In or after 221--204 BC or 145--117 BC; disputed alternatives}
\scripturenote{In \sourceurl{https://www.newadvent.org/cathen/15666a.htm}{``Book of Wisdom,'' \emph{Catholic Encyclopedia}}, Francis E. Gigot relates the book's persecution setting to either Ptolemy IV (221--204 BC) or Ptolemy VII (145--117 BC), allowing publication after their deaths. These are alternative lower bounds, not a closed date range. The verse recalls wilderness manna without establishing its event date.}
\scripturerow{Acclamation}{Matt.~4:4b}{The wilderness of Judea in Matthew's
narrative}{Narrated at the beginning of Jesus' public ministry}
\scripturenote{The liturgical line is an excerpt from Jesus' answer to the
tempter and itself cites Deuteronomy. Matthew's Gospel reached its early
Christian audience after the narrated event; the excerpt does not transfer
the temptation scene into John~6.}
\scripturerow{Gospel; Communion B}{John~6:\newline Gospel: 24--35;\newline Communion B: 35}{Across the Sea of Galilee;
the discourse begins at Capharnaum}{First-century narrated ministry; Gospel
composed later}
\scripturenote{John places the exchange on the day after the feeding and the
crossing. The crowd seeks Jesus, invokes manna, and hears the discourse move
from perishing food to the Father's present gift and Christ's
self-identification as bread of life. Communion B repeats v.~35.}
\scripturerow{Second Reading}{Eph.~4:17, 20--24}{Addressed to gentile
Christians in Asia Minor}{First century; Pauline authorship is traditional,
while modern scholarship disputes direct authorship}
\scripturenote{The selected verses belong to the letter's transition from
ecclesial confession to baptismal conduct. Verses 18--19 are not appointed;
the boundary therefore moves directly from the prohibition in v.~17 to
learning Christ in vv.~20--24.}
\end{scripturestable}
